\documentclass[11pt,a4paper]{article}
\usepackage[margin=1in]{geometry}
\usepackage{amsmath,amssymb,amsthm}
\usepackage{algorithm}
\usepackage{algpseudocode}
\usepackage{booktabs}
\usepackage{hyperref}

\theoremstyle{definition}
\newtheorem{definition}{\textbf{Definition}}[section]
\newtheorem{theorem}{\textbf{Theorem}}[section]
\newtheorem{lemma}[theorem]{\textbf{Lemma}}
\newtheorem{remark}{\textbf{Remark}}[section]

\title{\textbf{Neural Ordinary Differential Equations}}
\author{\textbf{Team Scaibu} \\ \small Email: \texttt{jaydeep.9664920749@gmail.com} \\ \small \textbf{Architecture and Core Engine Team}}
\date{\textbf{October 2026}}

\begin{document}
\maketitle

\section{Definitions and Cognitive Mental Models}

\subsection{Formal Mathematical Definition}
\textbf{Definition (Continuous-Depth Dynamics and Adjoint Sensitivity State):}
Let $\mathbf{h}(t) \in \mathbb{R}^D$ denote a continuous hidden state trajectory parameterized by continuous depth coordinate $t \in [t_0, t_1]$. A \textbf{Neural Ordinary Differential Equation (Neural ODE)} replaces discrete residual layer transitions with a continuous vector field parameterized by neural network $f(\mathbf{h}(t), t; \theta)$:
{\boldmath
\begin{equation}
\frac{\mathrm{d}\mathbf{h}(t)}{\mathrm{d}t} = f(\mathbf{h}(t), t; \theta), \quad \mathbf{h}(t_1) = \mathbf{h}(t_0) + \int_{t_0}^{t_1} f(\mathbf{h}(t), t; \theta) \mathrm{d}t = \operatorname{ODESolve}(\mathbf{h}(t_0), f, t_0, t_1, \theta)
\end{equation}
}
Under the classical \textbf{Fourth-Order Runge-Kutta (RK4)} integration scheme with step size $\Delta t = t_1 - t_0$, the solution is computed algebraically via four intermediate stage slopes:
{\boldmath
\begin{align}
\mathbf{k}_1 &= f(\mathbf{h}(t_0), t_0; \theta) \\
\mathbf{k}_2 &= f\left( \mathbf{h}(t_0) + \frac{\Delta t}{2}\mathbf{k}_1, \, t_0 + \frac{\Delta t}{2}; \, \theta \right) \\
\mathbf{k}_3 &= f\left( \mathbf{h}(t_0) + \frac{\Delta t}{2}\mathbf{k}_2, \, t_0 + \frac{\Delta t}{2}; \, \theta \right) \\
\mathbf{k}_4 &= f\left( \mathbf{h}(t_0) + \Delta t \, \mathbf{k}_3, \, t_0 + \Delta t; \, \theta \right) \\
\mathbf{h}(t_1) &= \mathbf{h}(t_0) + \frac{\Delta t}{6} \left( \mathbf{k}_1 + 2\mathbf{k}_2 + 2\mathbf{k}_3 + \mathbf{k}_4 \right)
\end{align}
}
To train $\theta$ in $\mathcal{O}(1)$ memory without caching intermediate forward activations, the \textbf{Adjoint Sensitivity Method} defines the adjoint state vector $\mathbf{a}(t) = \frac{\partial \mathcal{L}}{\partial \mathbf{h}(t)}$, which satisfies the reverse-time continuous adjoint ODE:
{\boldmath
\begin{equation}
\frac{\mathrm{d}\mathbf{a}(t)}{\mathrm{d}t} = -\mathbf{a}(t)^T \frac{\partial f(\mathbf{h}(t), t; \theta)}{\partial \mathbf{h}}, \quad \frac{\mathrm{d}\mathcal{L}}{\mathrm{d}\theta} = -\int_{t_1}^{t_0} \mathbf{a}(t)^T \frac{\partial f(\mathbf{h}(t), t; \theta)}{\partial \theta} \mathrm{d}t
\end{equation}
}

\subsection{Expanded Non-Technical Definition (Intuition for Anyone)}
\textbf{Intuitive Conceptual Definition:}
\begin{enumerate}
    \item \textbf{Replacing Staircases with a Smooth Slide}: Standard deep learning models (like ResNets) are staircases with 100 discrete steps. Each step computes $h_{t+1} = h_t + f(h_t)$. Neural ODEs replace the staircase with a perfectly smooth slide: the depth is continuous time $t \in [0, 1]$.
    \item \textbf{Adaptive Precision (Thinking Harder when Needed)}: In a standard model, an easy input (a clear photo of a dog) takes the exact same 100 steps as a difficult, blurry photo. An ODE solver automatically adapts its precision: it glides through easy problems in 5 quick jumps, and takes 200 tiny micro-steps only when navigating sharp, turbulent bends.
    \item \textbf{The Time-Travel Trick (Constant Memory)}: Training standard 1,000-layer models requires gigabytes of memory to store every layer's activations for backpropagation. The Adjoint Method runs the ODE backwards in time during the backward pass! You reconstruct the entire forward movie backwards from the end, training infinitely deep networks in constant $\mathcal{O}(1)$ RAM!
    \item \textbf{Irregular Timestamps}: Perfect for messy real-world time-series (e.g. hospital patients whose heart rate is checked irregularly every 3 hours or 45 minutes), because an ODE can be evaluated at any fractional continuous timestamp.
\end{enumerate}

\subsection{Child-Friendly Mental Model (Physical Metaphor)}
Imagine dropping a leaf into a winding river:
\begin{itemize}
    \item \textbf{Discrete ResNet}: You take a photo of the leaf every 10 meters, guessing where it will land next.
    \item \textbf{Neural ODE}: You write down the mathematical formula for the water currents. If the river is calm and straight, you calculate where the leaf ends up 1 kilometer away in 1 second. If there are whirlpools, you slow down and trace the smooth swirling path without ever losing track!
\end{itemize}

\section{Core Mathematical Equations and Definitions}

\subsection{Continuous State Evolution Law}
{\boldmath
\begin{equation}
\mathbf{h}(t_1) = \mathbf{h}(t_0) + \int_{t_0}^{t_1} f(\mathbf{h}(\tau), \tau; \theta) \mathrm{d}\tau
\end{equation}
}
\begin{itemize}
    \item \textbf{Formal Definition}: The integral formulation of continuous hidden state dynamics across integration interval $[t_0, t_1]$.
    \item \textbf{What It Means}: State transformation corresponds to accumulating an infinitesimal neural velocity vector field.
    \item \textbf{Why It Is Important}: Generalizes discrete layer depth to an uncountable continuum of infinitesimal transformations.
\end{itemize}

\subsection{Classical Runge-Kutta (RK4) Quadrature Step}
{\boldmath
\begin{equation}
\mathbf{h}_{t+\Delta t} = \mathbf{h}_t + \frac{\Delta t}{6} \left( \mathbf{k}_1 + 2\mathbf{k}_2 + 2\mathbf{k}_3 + \mathbf{k}_4 \right)
\end{equation}
}
\begin{itemize}
    \item \textbf{Formal Definition}: The explicit 4-stage Runge-Kutta numerical integration scheme with Simpson's rule quadrature weighting.
    \item \textbf{What It Means}: Samples the velocity field at the beginning, two midpoint estimates, and the endpoint, canceling Taylor terms up to order 4.
    \item \textbf{Why It Is Important}: Exhibits $\mathcal{O}(\Delta t^5)$ local truncation error, providing high-precision trajectory resolution with fixed step size.
\end{itemize}

\subsection{Continuous Backward Adjoint State Differential Equation}
{\boldmath
\begin{equation}
\frac{\mathrm{d}\mathbf{a}(t)}{\mathrm{d}t} = -\mathbf{a}(t)^T \mathbf{J}_f(\mathbf{h}(t), t) = -\sum_{j=1}^D a_j(t) \frac{\partial f_j(\mathbf{h}(t), t)}{\partial \mathbf{h}}
\end{equation}
}
\begin{itemize}
    \item \textbf{Formal Definition}: The differential equation governing the instantaneous sensitivity vector $\mathbf{a}(t) = \frac{\partial \mathcal{L}}{\partial \mathbf{h}(t)}$ running backward from $t_1$ to $t_0$.
    \item \textbf{What It Means}: Transports upstream loss gradients backward along the state trajectory via vector-Jacobian products.
    \item \textbf{Why It Is Important}: Allows evaluating exact parameter gradients by solving an augmented ODE backward without caching forward hidden states.
\end{itemize}

\subsection{Augmented State Dynamics for Parameter Gradients}
{\boldmath
\begin{equation}
\frac{\mathrm{d}}{\mathrm{d}t} \begin{bmatrix} \mathbf{h}(t) \\ \mathbf{a}(t) \\ \frac{\partial \mathcal{L}}{\partial \theta} \end{bmatrix} = \begin{bmatrix} f(\mathbf{h}(t), t; \theta) \\ -\mathbf{a}(t)^T \frac{\partial f}{\partial \mathbf{h}} \\ -\mathbf{a}(t)^T \frac{\partial f}{\partial \theta} \end{bmatrix}
\end{equation}
}
\begin{itemize}
    \item \textbf{Formal Definition}: The composite system of ordinary differential equations simultaneously integrating states, adjoints, and parameter gradients.
    \item \textbf{What It Means}: Concatenates the forward state, backward gradient, and parameter sensitivity into a single unified backward ODE solve.
    \item \textbf{Why It Is Important}: Guarantees constant $\mathcal{O}(1)$ memory scaling with respect to integration depth.
\end{itemize}

\section{Rigorous Mathematical Analysis Checklist}

\subsection{Notation and Symbol Semantics}
\begin{itemize}
    \item $\mathbf{h}(t) \in \mathbb{R}^D$: Continuous hidden state at time $t$.
    \item $f(\mathbf{h}, t; \theta) \in \mathbb{R}^D$: Parameterized velocity neural network.
    \item $\mathbf{k}_1, \mathbf{k}_2, \mathbf{k}_3, \mathbf{k}_4 \in \mathbb{R}^D$: Intermediate RK4 derivative slopes.
    \item $\Delta t \in \mathbb{R}^+$: Time step size.
    \item $\mathbf{a}(t) = \frac{\partial \mathcal{L}}{\partial \mathbf{h}(t)} \in \mathbb{R}^D$: Adjoint sensitivity state vector.
\end{itemize}

\subsection{Epistemic Status Table}
\begin{table}[h]
\centering
\small
\begin{tabular}{@{}llll@{}}
\toprule
\textbf{Quantity} & \textbf{Symbol} & \textbf{Epistemic Status} & \textbf{Operational Role} \\
\midrule
Continuous State & $\mathbf{h}(t)$ & Dynamic Solution Variable & Carries representation along continuous depth \\
Vector Field & $f(\mathbf{h}, t)$ & Neural Function Approximation & Defines local trajectory tangent \\
Effective Drift & $\bar{\mathbf{k}}$ & Simpson Quadrature Average & Discretized RK4 integration step direction \\
Adjoint State & $\mathbf{a}(t)$ & Continuous Lagrange Multiplier & Propagates gradients backward in time \\
\bottomrule
\end{tabular}
\end{table}

\subsection{Computational Dependency Graph}
{\centering
$\mathbf{h}_0 \longrightarrow \{\mathbf{k}_1, \mathbf{k}_2, \mathbf{k}_3, \mathbf{k}_4\} \longrightarrow \mathbf{h}_1 \longrightarrow \mathcal{L} \longrightarrow \mathbf{a}_1 \longrightarrow \text{Adjoint ODE} \longrightarrow \{\mathbf{a}_0, \nabla_\theta \mathcal{L}\}$
\par}

\subsection{Dimension and Coordinate Consistency}
All states $\mathbf{h}(t)$, slopes $\mathbf{k}_i$, drifts, and adjoints $\mathbf{a}(t)$ reside in $\mathbb{R}^D$. Continuous time $t$ and step $\Delta t$ reside in $\mathbb{R}$.

\subsection{Asymptotic Regimes}
\begin{itemize}
    \item \textbf{As $\Delta t \to 0$}: Discrete RK4 trajectory converges to the analytical Picard solution with global error $\mathcal{O}(\Delta t^4)$.
    \item \textbf{As $t_1 \to \infty$}: Trajectory approaches continuous dynamical system attractors or limit cycles.
\end{itemize}

\subsection{Boundary Constraints and Invariants}
\begin{itemize}
    \item \textbf{Reversibility Invariant}: If $f$ is Lipschitz continuous, the trajectory $\mathbf{h}(t)$ is uniquely reversible forward and backward in time.
\end{itemize}

\subsection{Structural Isomorphisms}
Neural ODEs are structurally isomorphic to Pontryagin's Minimum Principle in optimal control theory and continuous-time Hamiltonian mechanics.

\section{Theoretical Theorems and Formal Proofs}

\begin{theorem}[Pontryagin's Adjoint State Dynamics for Neural ODEs]
Let $\mathcal{L}(\mathbf{h}(t_1))$ be a scalar loss function evaluated on terminal state $\mathbf{h}(t_1)$, where $\mathbf{h}(t)$ evolves via $\frac{\mathrm{d}\mathbf{h}}{\mathrm{d}t} = f(\mathbf{h}(t), t; \theta)$. Then the adjoint state $\mathbf{a}(t) = \frac{\partial \mathcal{L}}{\partial \mathbf{h}(t)}$ satisfies the continuous differential equation:
{\boldmath
\begin{equation}
\frac{\mathrm{d}\mathbf{a}(t)}{\mathrm{d}t} = -\mathbf{a}(t)^T \frac{\partial f(\mathbf{h}(t), t; \theta)}{\partial \mathbf{h}(t)}
\end{equation}
}
with terminal boundary condition $\mathbf{a}(t_1) = \frac{\partial \mathcal{L}}{\partial \mathbf{h}(t_1)}$.
\end{theorem}

\begin{proof}
Let $t$ and $t + \epsilon$ be two points along the trajectory separated by an infinitesimal interval $\epsilon > 0$.
By first-order Euler expansion:
{\boldmath
\begin{equation}
\mathbf{h}(t + \epsilon) = \mathbf{h}(t) + \epsilon f(\mathbf{h}(t), t; \theta) + \mathcal{O}(\epsilon^2)
\end{equation}
}
Applying the multivariable chain rule to the gradient $\mathbf{a}(t) = \frac{\partial \mathcal{L}}{\partial \mathbf{h}(t)}$:
{\boldmath
\begin{equation}
\mathbf{a}(t) = \frac{\partial \mathcal{L}}{\partial \mathbf{h}(t + \epsilon)} \frac{\partial \mathbf{h}(t + \epsilon)}{\partial \mathbf{h}(t)} = \mathbf{a}(t + \epsilon) \frac{\partial}{\partial \mathbf{h}(t)}\left[ \mathbf{h}(t) + \epsilon f(\mathbf{h}(t), t; \theta) + \mathcal{O}(\epsilon^2) \right]
\end{equation}
}
Evaluating the Jacobian matrix:
{\boldmath
\begin{equation}
\frac{\partial \mathbf{h}(t + \epsilon)}{\partial \mathbf{h}(t)} = \mathbf{I} + \epsilon \frac{\partial f(\mathbf{h}(t), t; \theta)}{\partial \mathbf{h}(t)} + \mathcal{O}(\epsilon^2)
\end{equation}
}
Multiplying by row vector $\mathbf{a}(t + \epsilon)$:
{\boldmath
\begin{equation}
\mathbf{a}(t) = \mathbf{a}(t + \epsilon) + \epsilon \, \mathbf{a}(t + \epsilon) \frac{\partial f(\mathbf{h}(t), t; \theta)}{\partial \mathbf{h}(t)} + \mathcal{O}(\epsilon^2)
\end{equation}
}
Rearranging terms to construct the derivative definition:
{\boldmath
\begin{equation}
\frac{\mathbf{a}(t + \epsilon) - \mathbf{a}(t)}{\epsilon} = -\mathbf{a}(t + \epsilon) \frac{\partial f(\mathbf{h}(t), t; \theta)}{\partial \mathbf{h}(t)} + \mathcal{O}(\epsilon)
\end{equation}
}
Taking the limit as $\epsilon \to 0$:
{\boldmath
\begin{equation}
\frac{\mathrm{d}\mathbf{a}(t)}{\mathrm{d}t} = \lim_{\epsilon \to 0} \frac{\mathbf{a}(t + \epsilon) - \mathbf{a}(t)}{\epsilon} = -\mathbf{a}(t)^T \frac{\partial f(\mathbf{h}(t), t; \theta)}{\partial \mathbf{h}(t)}
\end{equation}
}
This establishes the continuous adjoint differential equation, with terminal boundary condition $\mathbf{a}(t_1) = \nabla_{\mathbf{h}(t_1)} \mathcal{L}$ given directly by the loss gradient.
\end{proof}

\begin{theorem}[Order of Local and Global Error for RK4 Integration]
Assume $f(\mathbf{h}, t)$ is four times continuously differentiable ($\mathcal{C}^4$) with Lipschitz constant $L$. The local truncation error of the classical RK4 step is $\mathcal{O}(\Delta t^5)$, and the accumulated global error over interval $[t_0, t_1]$ is $\mathcal{O}(\Delta t^4)$.
\end{theorem}

\begin{proof}
Let $\mathbf{h}(t + \Delta t)$ be the exact solution. Expanding in a Taylor series about $t$:
{\boldmath
\begin{equation}
\mathbf{h}(t + \Delta t) = \mathbf{h}(t) + \Delta t \mathbf{h}' + \frac{\Delta t^2}{2} \mathbf{h}'' + \frac{\Delta t^3}{6} \mathbf{h}''' + \frac{\Delta t^4}{24} \mathbf{h}^{(4)} + \mathcal{O}(\Delta t^5)
\end{equation}
}
Expanding each RK4 evaluation $\mathbf{k}_1, \mathbf{k}_2, \mathbf{k}_3, \mathbf{k}_4$ via multivariable Taylor expansions in $(\mathbf{h}, t)$ and substituting into the weighted sum:
{\boldmath
\begin{equation}
\mathbf{h}_{\text{RK4}} = \mathbf{h}(t) + \frac{\Delta t}{6}(\mathbf{k}_1 + 2\mathbf{k}_2 + 2\mathbf{k}_3 + \mathbf{k}_4)
\end{equation}
}
By matching algebraic coefficients with the exact Taylor expansion, the terms of order $\Delta t, \Delta t^2, \Delta t^3,$ and $\Delta t^4$ match identically:
{\boldmath
\begin{equation}
\|\mathbf{h}(t + \Delta t) - \mathbf{h}_{\text{RK4}}\|_2 = \mathcal{O}(\Delta t^5)
\end{equation}
}
Summing the local errors over $N = (t_1 - t_0)/\Delta t$ steps, by Gronwall's lemma the accumulated global error satisfies:
{\boldmath
\begin{equation}
E_{\text{global}} \le \frac{C}{L} (e^{L(t_1 - t_0)} - 1) \Delta t^4 = \mathcal{O}(\Delta t^4)
\end{equation}
}
Thus, halving the integration step size $\Delta t$ reduces the global numerical error by $2^4 = 16\times$.
\end{proof}

\section{Foundational Architectural Questions and Epistemic Meta-Questions}

\subsection{Architectural Question 1: Constant Memory vs Reverse Numerical Instability}
\textbf{Question:} Does the adjoint sensitivity method always eliminate memory bottlenecks in practical training?\\
\textbf{Answer:} While the adjoint method achieves $\mathcal{O}(1)$ memory scaling, integrating stiff or chaotic ODEs backward in time can be numerically unstable due to accumulated solver discretization errors. For stiff dynamics, Checkpointing Adjoint methods (e.g. Interpolated Reverse-Time ODE) cache a few checkpoints along the trajectory to ensure numerical stability.

\textbf{Meta-Question:} How can one detect if backward reconstruction has suffered numerical drift?\\
\textbf{Answer:} By computing the discrepancy $\|\mathbf{h}_{\text{reconstructed}}(t_0) - \mathbf{h}_{\text{true}}(t_0)\|_2$. If the reconstruction error exceeds a prescribed tolerance, the solver step size $\Delta t$ must be shrunk or checkpointing activated.

\subsection{Architectural Question 2: Function Evaluation Growth (NFE Explosion)}
\textbf{Question:} Why does the number of function evaluations (NFE) often increase dramatically during Neural ODE training?\\
\textbf{Answer:} Unregularized neural networks often learn vector fields with high spatial frequency and extreme curvature. High curvature forces adaptive step-size ODE solvers to take increasingly tiny steps to maintain local tolerance, causing training epochs to slow down by $10\times$. Regularizing the Frobenius norm of the Jacobian $\mathbf{J}_f$ or penalizing kinetic energy $\int \|f\|^2 \mathrm{d}t$ keeps NFE low and constant.

\textbf{Meta-Question:} Can Neural ODEs represent intersecting trajectories?\\
\textbf{Answer:} No. By the Picard-Lindel\"of theorem, solutions to first-order ODEs cannot intersect in $\mathbb{R}^D$. If a task requires trajectories to cross (e.g. concentric circles in 1D), Augmented Neural ODEs add dummy zero-initialized dimensions $\mathbb{R}^{D + P}$ to allow trajectories to lift and cross in higher-dimensional space without intersecting.

\section{Algorithmic Implementation Specification}

\begin{algorithm}[H]
\caption{Fourth-Order Runge-Kutta (RK4) Discrete ODE Integration Step}
\begin{algorithmic}[1]
\Require State $\mathbf{h} \in \mathbb{R}^D$, slopes $\mathbf{k}_1, \mathbf{k}_2, \mathbf{k}_3, \mathbf{k}_4 \in \mathbb{R}^D$, step size $\Delta t > 0$
\Ensure Integrated state $\mathbf{h}_{\text{next}} \in \mathbb{R}^D$, effective drift $\bar{\mathbf{k}} \in \mathbb{R}^D$, displacement norm $\|\Delta \mathbf{h}\|_2$
\State Initialize empty arrays $\mathbf{h}_{\text{next}}, \bar{\mathbf{k}}$ of length $D$
\State $disp\_sq \gets 0.0, \quad inv\_six \gets 1.0 / 6.0$
\For{$d = 0$ \textbf{to} $D - 1$}
    \State $drift \gets inv\_six \cdot (\mathbf{k}_1[d] + 2.0 \cdot \mathbf{k}_2[d] + 2.0 \cdot \mathbf{k}_3[d] + \mathbf{k}_4[d])$
    \State $disp \gets \Delta t \cdot drift$
    \State $disp\_sq \gets disp\_sq + disp \cdot disp$
    \State $\bar{\mathbf{k}}[d] \gets drift$
    \State $\mathbf{h}_{\text{next}}[d] \gets \mathbf{h}[d] + disp$
\EndFor
\State \Return $(\mathbf{h}_{\text{next}}, \bar{\mathbf{k}}, \sqrt{disp\_sq})$
\end{algorithmic}
\end{algorithm}

\section{Big-$\Theta$ Complexity Analysis}

\begin{itemize}
    \item \textbf{Integration Step Time Complexity}: $\Theta(D)$ vector arithmetic operations per RK4 stage.
    \item \textbf{Total Trajectory Time Complexity}: For $N_{\text{steps}}$ solver intervals, total runtime is $\Theta(N_{\text{steps}} \cdot (D + 4 \mathcal{C}_{\text{NN}}))$.
    \item \textbf{Space Complexity}: $\Theta(D)$ working buffers for state and slope vectors.
    \item \textbf{Work \& Depth}: Work is $\mathcal{W} = \Theta(N_{\text{steps}} \cdot D)$; parallel depth across coordinates is $\mathcal{D} = \Theta(\log D)$.
\end{itemize}

\section{Single Canonical Repository Reference}

The complete deterministic scalar implementation, verification suite, and unit test benchmarks are published at:
\begin{center}
\url{https://github.com/Chief-Strategist-J/llm-obs-infra}
\end{center}

\end{document}
